\documentclass[11pt]{article}
\usepackage[T1]{fontenc}
\usepackage[includeheadfoot,margin=0.8in,top=0.4in,bottom=0.2in]{geometry}
\usepackage[hidelinks]{hyperref}
\usepackage{lastpage}
\usepackage{mathpazo}
\usepackage{setspace}
\usepackage{fancyhdr}
\pagestyle{fancy}
\fancyhf{}
\fancyhead[L]{\small\textbf{Yumeng He}}
\fancyhead[R]{\small Web: \href{https://heyumeng.com}{heyumeng.com} \quad Email: \href{mailto:heyumeng0928@gmail.com}{heyumeng0928@gmail.com}}
\renewcommand{\headrulewidth}{0.2pt}
\setlength{\headheight}{16pt}
\setlength{\headsep}{6pt}
\fancyfoot[C]{\small Page \thepage\ of \pageref{LastPage}}
\setlength{\footskip}{18pt}

% Slightly increase line spacing across the document
\setstretch{1.0}

% Add slight space between paragraphs (adjust here)
\setlength{\parskip}{0.3em}

\begin{document}

\vspace*{-18pt}
\begingroup
\centering
\Large\bfseries Cyberbullying Awareness and Prevention \par
\endgroup
% \vspace{6pt}

What would you do if a naked photo of your friend appeared on your phone? Or a
sexual video that seemed to capture their voice and gestures perfectly? Telling
yourself not to share it is only the beginning. In 2025, the speed and cruelty
of cyberbullying have grown alongside tools that make images and videos look
impossibly real. I write this as a master student in computer science who works
on computer graphics and vision, especially 2D and 3D generation. 
I once thought my field was only about creativity and
beauty. Now I know it is also a double-edged sword. With a single sentence, a face can be
conjured. With one image, a system can fabricate ten more that appear to show
the same person in places they never were. The technology is not the villain.
Like a knife, it can slice fruit or draw blood. 
The difference is what people choose to do and what communities are prepared to prevent.

Parents hold the first line. When a harmful image spreads, minutes matter.
Shame hardens quickly into silence. Families can rehearse the moment before it
happens. During a car ride or at dinner, a parent can ask what a child would do
if they received an abusive screenshot, then practice a response. The plan can
be simple. Pause. Do not forward. Save enough context to preserve evidence. Tell
a trusted adult and the person targeted right away. When a child admits a mistake, the first response should be
safety and care, not punishment. That keeps the door open the next time. 

School leaders have a different kind of power. They can build systems that move
faster than rumor. Every school should have a plain language policy that defines
cyberbullying, nonconsensual imagery, impersonation, and deepfake harassment.
Students and staff should know exactly how to report, where to go, and what will
happen next. An anonymous channel lowers the barrier. A rapid response protocol
can state who contacts families, who preserves evidence with care, and who
checks on the student between classes. Consequences can be firm, but the aim
should be learning and repair rather than spectacle. 

As for me, my role is to push technology forward while building the guardrails
that keep people safe. I design and advocate for safety by default. That means
supporting invisible tags that travel with media and quietly declare how it was
made. Watermarking, content credentials, and tamper evident signatures can give
images and videos a kind of birth certificate. When a photo is posted, newsrooms
and platforms can read those signals and show users the provenance. I can help
test these methods in my lab, measuring how robust they are when files are
cropped, resized, or repeatedly shared. I can publish results that are honest
about tradeoffs so the field learns where these tools work and where they need
strengthening. In my code releases, I can include risk analysis, usage
guidelines, and clear options that enable provenance by default. I can choose
training data and model settings that reduce the chance of impersonation and
abuse, and I can limit public demos that would invite misuse. I can partner with
educators, counselors, and journalists to bring verification rituals into daily
practice. Reverse image search. Ask for the source and the date. Look for
reflections, edges, and textures that repeat. Most of all, I can help build a
culture where astonishing images are treated as claims that must be proven, and
where the easiest path for a platform is to surface the history of an image
rather than simply chase clicks.

When the worst happens, our reactions should be both practical and humane. We
can keep copies of the evidence somewhere safe, using only what is needed to
report to the platform or the school. We can comfort without asking the person
targeted to retell the story again and again. We can remind them that the shame
belongs to the people who manufactured harm, not to them. 

Prevention is more than a list of rules. It is an atmosphere. It is the feeling
that if something goes wrong online, you will be believed, you will be helped,
and people will act. Parents grow that feeling at home. School leaders build it
into policies that are clear and humane. In the lab, I try to steer new tools
toward accountability by giving media a history that can be read and trusted.
Technology will keep advancing. Our compassion and our preparedness must advance
with it, so that the images we share do not only carry pixels but also the truth
of how they were made.
\end{document}
